\subsection{切丛}

8.1.1. 设 \(r \in \mathbf{N}_{\mathbb{K}}\)，X 是类为 \(\mathbf{C}^r\) 的 K-流形。记 T(X) 为切空间 \(\mathrm{T}_x(\mathrm{X})\)（5.5.1）的不交并，其中 \(x \in \mathrm{X}\)，并记从 T(X) 到 X 的典范映射为 \(\pi\)。设 \(c = (\mathrm{U}, \varphi, \mathrm{E})\) 是流形 X 的一个图；对于 \(x \in \mathrm{U}\) 和 \(t \in \mathrm{T}_x(\mathrm{X})\)，置 \(\zeta_c(x, t) = (x, \theta_c^{-1}(t))\)，其中映射 \(\theta_c\) 是 5.5.1 中定义的映射。三元组 \(c' = (\mathrm{U}, \zeta_c, \mathrm{E})\) 是 T(X) 的一个向量丛图（7.1.1）（T(X) 赋予映射 \(\pi: \mathrm{T}(\mathrm{X}) \to \mathrm{X}\)）。在 T(X) 上存在一个以 X 为基底、类为 \(\mathbf{C}^{r-1}\) 的向量丛结构，唯一地使得对于 X 的每个图 \(c\)，相应的向量丛图 \(c'\) 是 T(X) 的一个向量丛图 \(^{(1)}\)。赋予此结构后，T(X) 称为 X 的切丛。特别地，T(X) 具有类为 \(\mathbf{C}^{r-1}\) 的流形结构，且三元组（T(X)，X，\(\pi\)）是一个纤维化（7.1.4）。

设 U 是巴拿赫空间 E 的一个开子集，i 是 U 到 E 的典范嵌入。由向量丛图 \(c' = (U, \zeta_c, E)\)（它与来自 U 的图 \(c = (U, i, E)\) 相关联）定义了从 T(U) 到平凡丛 \(E_U\) 的一个同构（7.1.5）；我们借助此同构将这两个丛视为相同。

若 X 是 E 型纯流形（5.1.7），则 T(X) 是 E \(\times\) E 型纯流形。

8.1.2. 设 \(f: X \to Y\) 是类为 \(\mathbf{C}^r\) 的流形态射。我们定义一个 \(f\)-态射 \(\mathbf{T}(f): \mathbf{T}(\mathbf{X}) \to \mathbf{T}(\mathbf{Y})\)，它是类为 \(\mathbf{C}^{r-1}\) 的向量丛态射，并规定 \(\mathbf{T}(f)(x, t) = (f(x), \mathbf{T}_x(f)(t))\)，其中 \(x \in \mathbf{X}\) 且 \(t \in \mathbf{T}_x(\mathbf{X})\)。

若 \(f = \operatorname{Id}_{\mathbf{X}}\)，则有 \(\mathbf{T}(f) = \operatorname{Id}_{\mathbf{T}(\mathbf{X})}\)。对于 \(f: \mathbf{X} \to \mathbf{Y}\) 和 \(g: \mathbf{Y} \to \mathbf{Z}\)，有 \(\mathbf{T}(g \circ f) = \mathbf{T}(g) \circ \mathbf{T}(f)\)。

8.1.3. 设 \(f: \mathbf{X} \to \mathbf{Y}\) 是类为 \(\mathbf{C}^r\) 的流形态射，并令 \(f': \mathbf{T}(\mathbf{X}) \to f^* \mathbf{T}(\mathbf{Y})\)

\(^{1}\) 当 K = R 且 r = 1 时，T(X) 是一个拓扑向量丛（参见记号和约定）。

它是满足 \(\mathrm{T}(f) = g \circ f'\) 的唯一 X-态射，其中 \(g\) 是典范态射 \(f^{*}\mathrm{T}(\mathrm{Y}) \to \mathrm{T}(\mathrm{Y})\)（参见 7.2.4）。

对于 \(f\) 是浸入（分别地，浸没），其充要条件是 \(0 \to \mathrm{T}(\mathrm{X}) \xrightarrow{f'} f^* \mathrm{T}(\mathrm{Y})\)（分别地，\(\mathrm{T}(\mathrm{X}) \xrightarrow{f'} f^* \mathrm{T}(\mathrm{Y}) \to 0\)）是一个局部直正合序列（7.5.6）。如果 \(f\) 是浸入，\(f'\) 的余核丛称为 \(f\) 的法丛或横截丛。如果 \(f\) 是浸没，\(f'\) 的核丛称为 \(f\) 的纤维的切丛，或称为 X 在 Y 上的相对切丛，并记为 \(\mathrm{T}(\mathrm{X}/\mathrm{Y})\)。其纤维 \(\mathrm{T}_x(\mathrm{X}/\mathrm{Y})\)，在 \(x \in \mathrm{X}\) 处，是子流形在 \(x\) 处的切空间，该子流形为 \(f^{-1}(f(x))\)。序列

\[
0 \longrightarrow \mathrm{T} _ {x} (\mathrm{X} / \mathrm{Y}) \stackrel {i} {\longrightarrow} \mathrm{T} _ {x} (\mathrm{X}) \stackrel {\mathrm{T} _ {x} (f)} {\longrightarrow} \mathrm{T} _ {f (x)} (\mathrm{Y}) \longrightarrow 0
\]

其中 i 是典范嵌入，序列是正合的。

\(f\) 是平展的充要条件是 \(f'\) 是一个同构。

当 K 的特征为 0 时，f 是子浸入的充要条件是 \(f'\) 是局部直的。

8.1.4. 设 \(X_{1}\) 和 \(X_{2}\) 是两个流形。我们有 \(\mathrm{T}(\mathrm{X}_{1} \times \mathrm{X}_{2}) = p_{1}^{*}\mathrm{T}(\mathrm{X}_{1}) \oplus p_{2}^{*}\mathrm{T}(\mathrm{X}_{2})\)，其中 \(p_{1}\) 和 \(p_{2}\) 是典范投影（5.6.3）。

8.1.5. 设 \(n\) 是一个满足 \(\geqslant 0\) 的整数；\(\mathrm{T}(\mathbf{K}^{n})\)（\(\mathbf{K}^{n}\) 的切丛）与纤维为 \(\mathbf{K}^{n}\) 的平凡丛相同（8.1.1）。由 \(\mathbf{K}^{n}\) 的典范基的元素定义的常值截面构成 \(\mathrm{T}(\mathbf{K}^{n})\) 的一个标架。

若 X 是一个流形，且 \((u^{1},\ldots,u^{n})\) 是 X 的开集 U 上的一个坐标系（5.1.10），则前述标架通过结构传递定义了 U 上 T(X) 的一个标架；该标架称为由 \((u^{1},\ldots,u^{n})\) 定义的切标架，记为 \((\partial/\partial u^{1},\ldots,\partial/\partial u^{n})\)（参见 5.5.8）。

